
Foundation models trained with billions of parameters on massive datasets require extensive data curation---deduplication, perplexity-based filtering, quality scoring---at each training stage. The C4 pre-training corpus applies deduplication threshold 0.7-0.8 and perplexity cutoff ~1000. When practitioners curate instruction fine-tuning data (Dolly, Alpaca) or RLHF preference datasets, they typically re-optimize these same thresholds from scratch. This practice raises a fundamental question: which curation decisions are universal across training stages, and which require stage-specific tuning?

Existing research treats each training stage in isolation. DataComp optimizes filtering strategies for vision-language pre-training. Alpagasus refines instruction dataset quality through ChatGPT-based scoring. RLHF preference selection operates independently of earlier curation decisions. No systematic understanding exists of which curation components transfer robustly across the foundation model pipeline versus which depend critically on stage-specific objectives.

We hypothesize that curation operations vary in their dependence on training objectives. Low-level quality filters like deduplication and perplexity-based outlier removal operate on surface statistics---n-gram overlap, token probability distributions---that remain invariant to whether a model trains for next-token prediction (pre-training), instruction following (fine-tuning), or human preference alignment (RLHF). High-level strategies like domain mixing and task-specific filters directly encode assumptions about what the model should learn at each stage. This distinction predicts transfer behavior: operations with low mutual information I(filter\_decision; stage\_objective) should transfer robustly, while those tightly coupled to stage goals require re-tuning.

We test this hypothesis through four controlled experiments examining: (1) existence of transfer-stable curation categories, (2) cross-stage threshold invariance for low-level operations, (3) categorical separation between objective-independent and objective-dependent techniques, and (4) quality-speed trade-offs in embedding-based subset selection. Experiments use C4 pre-training data (52,002-sample subset), Dolly-15k and Alpaca-52k instruction fine-tuning datasets, and Llama-2-7B as the base model. Evaluation uses MMLU and HellaSwag benchmarks.

Our main findings are:

\textbf{Transfer stability taxonomy validated.} Low-level filters (deduplication, perplexity) transfer with 0.38\% average delta (MMLU and HellaSwag combined), while high-level strategies (instruction quality filters) show 5.04\% degradation when transferred. Non-overlapping 95\% confidence intervals ([0.30\%, 0.46\%] vs. [4.44\%, 5.65\%]) and large effect size (Cohen's d=10.76) confirm categorical separation. Welch's t-test yields p=0.078, slightly above conventional significance thresholds but consistent with small sample size (n=2 benchmarks per category).

\textbf{Optimal thresholds are stage-invariant.} Independent optimization for C4 pre-training versus Dolly fine-tuning yields identical thresholds (dedup=0.7, perplexity=500). Cross-stage threshold application incurs maximum 3.0\% performance penalty, well below the 10\% threshold defining transfer failure.

\textbf{Quality-speed trade-offs quantified.} Embedding-based subset selection using early-stage embeddings (MiniLM) achieves 6.$9\times$ speedup over late-stage embeddings (Instructor) with 2.4\% quality penalty at k=5000 subset size. Late-stage embeddings at k=10000 preserve 99.6\% of full-dataset quality (0.4\% degradation).

These contributions shift the paradigm from per-stage curation optimization to transfer-aware design. Practitioners can reuse C4 pre-training thresholds (dedup 0.7-0.8, perplexity 500-1000) for instruction fine-tuning without re-optimization, focusing tuning resources on domain mixing and task-specific filters.

\textbf{Methodological constraint.} All experiments were conducted at proof-of-concept (PoC) tier. Curation pipelines executed on full datasets, but model training was simulated using predicted scores derived from published benchmarks (Llama-2 technical report, DataComp studies). This validates mechanism direction---categorical separation, threshold invariance---but defers absolute precision to production execution. Relative patterns (transfer delta, effect sizes) are grounded in literature; production validation with full training and lm-evaluation-harness is needed to confirm absolute performance values.
